\section{A uniform estimate for several radius hits}
\label{mp:section}

The construction of a trace measure will require moments of every fixed
integer order. The necessary probability estimate allows each point to
reach its own target at a different time. Its constant is uniform over
the physical domain, a feature that will later be essential for disk
models with arbitrary boundary angles.

Multipoint Green estimates and their moment consequences have a substantial
history. Lawler--Werness \cite[Theorems~1--2]{LawlerWerness2013}
prove ordered two-point Green limits and unequal-radius approach estimates;
Rezaei--Zhan \cite[Theorems~1.1--1.2]{RezaeiZhan2017} prove all-order
Euclidean approach bounds and finite moments of natural content.
The estimate needed here is a consequence of those Euclidean bounds and
Koebe's estimates; a comparison after the proof gives the deduction.
We give a direct proof in arbitrary physical domains using regularized
products. Each completed point is frozen at its own hit, and a conformal
interaction potential controls the remaining factors. We then integrate
the collision kernel to obtain the required moments.

\begin{proposition}[Multipoint radius estimate]\label{mp:estimate}
Let $D=J(\HH)$ as in Section~\ref{op:section}, and let
$x_1,\ldots,x_m\in D$ be distinct. Suppose
\[
 0<b_i\le R_D(x_i),\qquad b_i\le |x_i-x_j|\quad(j\ne i).
\]
For targets $0<e_i\le b_i$, let $E$ be the event that each $x_i$
reaches conformal radius $e_i$ while alive. Then
\begin{equation}
 \PP(E)\le C_{m,a}\prod_{i=1}^m(e_i/b_i)^\alpha.
 \label{mp:hit-bound}
\end{equation}
The constant depends only on $m$ and $a$, independently of $D$, $J$,
the initial angles, the radii, and the order and times of the hits.
A point need not remain alive after reaching its own target.
For targets exceeding some caps, the corresponding constraints may be
dropped, giving, after enlarging $C_{m,a}$,
\begin{equation}
 \PP(E)\le C_{m,a}
                \prod_{i=1}^m\min\{1,(e_i/b_i)^\alpha\}.
 \label{mp:capped-hit-bound}
\end{equation}
\end{proposition}

\subsection{Controlling the interaction of two active points}

Let $\sigma_i$ and $T_i$ be the first times when point $i$ reaches
radius $b_i$ and $e_i$, respectively, while alive, with value infinity
if the corresponding radius is never reached. The point is active on
$[\sigma_i,T_i)$ while it remains alive. Put
\[
 U_i=\frac1{2a}\log\frac{b_i}{e_i},\qquad
 u_i(t)=\int_{\sigma_i}^t V_i(v)^2\,\dd v
\]
on its active interval. Thus $R_t(x_i)=b_i e^{-2au_i(t)}$, and it
finishes when $u_i=U_i$. We first take $e_i<b_i$; constraints with
$e_i=b_i$ may be dropped.

For two surviving points put
\[
 L_{ij}(t)=\log\left|
       \frac{Z_i(t)-\overline{Z_j(t)}}{Z_i(t)-Z_j(t)}\right|\ge0.
\]
The common Brownian driver cancels in the two differences. Consequently
\begin{equation}
 \frac{\dd}{\dd t}L_{ij}
 =\operatorname{Re}\left(-\frac a{Z_i\overline{Z_j}}
                            +\frac a{Z_iZ_j}\right)
 =-2aV_iV_j.
 \label{mp:interaction-derivative}
\end{equation}
This is also a conformal invariant: for a disk uniformization
$F:\DD\to D_t$ with $F(0)=x_i$ and $F(w)=x_j$, it equals
$-\log|w|$. Write $\rho=|w|$ and $L=-\log\rho$. The growth estimate
\eqref{op:koebe-growth} gives
\[
 \frac{|x_i-x_j|}{R_t(x_i)}
 \le\frac{\rho}{(1-\rho)^2}
 =\frac1{4\sinh^2(L/2)}.
\]
During joint activity, the separation assumption and
$R_t(x_i)=b_i e^{-2au_i}$ imply
\[
 L_{ij}\le2\operatorname{arsinh}\left(
       \frac12\sqrt{\frac{R_t(x_i)}{|x_i-x_j|}}\right)
 \le e^{-au_i}.
\]
Interchanging the two points proves
\begin{equation}
 L_{ij}(t)\le e^{-a\max(u_i(t),u_j(t))}
 \qquad\text{during joint activity}.
 \label{mp:interaction-decay}
\end{equation}
In particular, for $c\ge0$ with $2c<a$,
\begin{equation}
 \int_{\text{joint activity}}
       e^{c(u_i+u_j)}V_iV_j\,\dd t
 \le I(a,c):=\frac{e^{2c}}{2a(1-e^{2c-a})}.
 \label{mp:interaction-integral}
\end{equation}
To see this, joint activity is an interval, and
$v=\max(u_i,u_j)$ increases on it. On its portion with
$k\le v<k+1$, the total decrease available to $L_{ij}$ is at most
$e^{-ak}$, by \eqref{mp:interaction-decay}. Equation~\eqref{mp:interaction-derivative}
bounds the unweighted integral there by $e^{-ak}/(2a)$, while the
weight is at most $e^{2c(k+1)}$. Summing over $k\ge0$ proves
\eqref{mp:interaction-integral}, also for any truncated part of the
joint active interval. This argument is made in $D_t$ itself; no
bound on the derivative of the initial map $J$ enters it.

\subsection{Regularized factors and the product supermartingale}

The factor $\sin^p\theta$ in the one-point martingale has no positive
lower bound uniform in the angle at a target hit. We add a positive
regularizing term at both ends of each active interval, so activation
and completion have angle-uniform bounds. The added term decays away
from those ends: this keeps its accumulated drift bounded independently
of the interval length. The symmetric distance
$\min\{u,U_i-u\}$ records exactly that choice.

The exponent $q$ will supply negative angular drift. Young's Inequality
will introduce a regularization loss with exponent $M$, and we choose
the decay rate $\beta$ so that the pair interaction can absorb that loss.
Use the explicit values
\begin{equation}
 q=\tfrac12\min(1,p),\qquad
 \eta=\frac{1-q}{2-q}\in(0,\tfrac12),\qquad
 M=\frac{1+\eta}{1-2\eta}=\frac3q-2,\qquad
 \beta=\min\left(1,\frac a{4M}\right).
 \label{mp:parameters}
\end{equation}
Thus $0<q<\min(1,p)$ and $2M\beta\le a/2<a$. For an active point
define
\begin{equation}
 \delta_i(u)=e^{-\beta\min(u,U_i-u)},\qquad
 f_i(u,\theta)=\sin^p\theta+\delta_i(u)(1+\sin^q\theta),\qquad
 D_i=\frac{\delta_i}{f_i}\sin^{q-2}\theta_i.
 \label{mp:factors}
\end{equation}
The factor for point $i$ is $1$ before activation,
\[
 Q_i(t)=\frac13 e^{pu_i(t)/2}f_i(u_i(t),\theta_i(t))
\]
during activity, and the constant $e^{pU_i/2}/3$ after completion.
The activation jump is downward because $f_i(0,\theta)\le3$;
the completion jump is downward because $f_i(U_i,\theta)\ge1$.
After completion the point's factor is frozen, including if that point
is swallowed later. Figure~\ref{fig:individual-stops} illustrates the
separate endpoints.
\begin{figure}[t]
\centering
\begin{tikzpicture}[x=0.85cm,y=0.72cm,>=Stealth,font=\small]
  \draw[->] (0,1) -- (8.2,1) node[right] {$t$};
  \foreach \y/\i/\start/\finish in {0/1/1/4.2,-1.4/2/2.5/6.5,-2.8/3/0.6/5.2} {
    \node[left] at (0,\y) {$x_{\i}$};
    \draw[gray!60] (0,\y)--(\start,\y);
    \draw[very thick,blue!65!black] (\start,\y)--(\finish,\y);
    \draw[dashed,thick] (\finish,\y)--(7.6,\y);
    \fill (\start,\y) circle(1.5pt);
    \fill (\finish,\y) circle(1.5pt);
    \node[above] at (\start,\y) {$\sigma_{\i}$};
    \node[above] at (\finish,\y) {$T_{\i}$};
  }
  \node[align=center] at (2.2,-4) {solid: active factor};
  \node[align=center] at (6,-4) {dashed: frozen factor};
\end{tikzpicture}
\caption{The multipoint event uses each point's own activation and target
stop. The solid intervals may overlap in any order. After $T_i$ the
factor for $x_i$ is constant, and no further survival of that point is
required. The axis is schematic physical time, not the individual
conformal-radius clocks.}
\label{fig:individual-stops}
\end{figure}


Put $Q=\prod_{i=1}^m Q_i$. If all points finish, this product has
the deterministic value
\[
 Q=3^{-m}\prod_i(b_i/e_i)^\alpha.
\]
Thus a uniform bound on the expectation of a discounted product will
give the desired probability estimate, provided the total discount is
deterministically bounded. The individual negative drifts absorb
the singular angular terms, and the pair interaction estimate pays for
what remains. We give both estimates explicitly. Write
$s=\sin\theta$, $k=q(p-q)/2$, and $\sigma=\delta'/\delta$, so
$|\sigma|\le\beta\le1$ almost everywhere. Formula~\eqref{op:sine-generator}
gives
\begin{equation}
 \frac{(\partial_u+\mathcal L+p/2)f}{f}
 =\frac\delta f\left[
 \frac p2+\frac{(p-q)(1+q)}2s^q-ks^{q-2}
                         +\sigma(1+s^q)\right].
 \label{mp:exact-drift}
\end{equation}
Let
\[
 A=\frac p2+\frac{(p-q)(1+q)}2+2,
 \qquad h_0=\left(\frac{k}{2A}\right)^{1/(2-q)}\in(0,1).
\]
The bracket in \eqref{mp:exact-drift} is at most $A-ks^{q-2}$.
For $s<h_0$, adding $(k/2)D$ leaves a nonpositive expression.
For $s\ge h_0$, use $f\ge s^p\ge h_0^p$. We obtain
\begin{equation}
 \frac{(\partial_u+\mathcal L+p/2)f}{f}
 \le C_1\delta-c_1D,
 \qquad C_1=Ah_0^{-p},\quad c_1=k/2>0.
 \label{mp:drift-bound}
\end{equation}
All these constants depend only on $a$. The corner of $\delta$ at
$U_i/2$ causes no additional term: it is an absolutely continuous
function of the finite-variation clock $u_i$.

Because $q<p$, $q<1$, and $0<\delta\le1$, angular differentiation
gives $|\partial_\theta f|\le(p+q)s^{q-1}$. Also $\delta\le f\le3$,
and $\eta(q-2)=q-1$. Therefore
\begin{equation}
 \frac{|\partial_\theta f|}{f}
 \le(p+q)\delta^{-1}s^{q-1}
 \le(p+q)3^\eta\delta^{-1-\eta}D^\eta.
 \label{mp:derivative-bound}
\end{equation}
This bound includes the singular angular behavior when $p<1$.

On active intervals, the normalized
Brownian coefficient of $Q_i$ is
$(\partial_\theta f_i/f_i)V_i$. It\^o's Product Formula thus gives
one cross term for each unordered active pair. Its absolute value is
bounded, by \eqref{mp:derivative-bound}, by
\[
 C_2(\delta_i\delta_j)^{-1-\eta}
    (D_iV_i^2)^\eta(D_jV_j^2)^\eta(V_iV_j)^{1-2\eta}.
\]
Weighted Young's Inequality with weights $\eta,\eta,1-2\eta$ shows
that for any $\epsilon>0$ this is at most
\begin{equation}
 \epsilon(D_iV_i^2+D_jV_j^2)
       +C_\epsilon(\delta_i\delta_j)^{-M}V_iV_j.
 \label{mp:young}
\end{equation}
Indeed, putting $r=1-2\eta$, the weighted arithmetic--geometric
mean inequality gives the displayed bound with
$C_\epsilon=rC_2^{1/r}(\eta/\epsilon)^{2\eta/r}$, and raises
$\delta_i\delta_j$ to exponent $-(1+\eta)/r=-M$.
For $m>1$ choose $\epsilon=c_1/[2(m-1)]$; for $m=1$ there are no
cross terms. Summing \eqref{mp:drift-bound} and \eqref{mp:young}
absorbs the pair contributions into the negative terms. A nonnegative
upper bound for the remaining normalized physical-time drift is
\[
 k_t=C_1\sum_{i\text{ active}}\delta_iV_i^2
       +C_3\sum_{\substack{i<j\\i,j\text{ active}}}
                     (\delta_i\delta_j)^{-M}V_iV_j,
\]
where $C_3$ depends only on $m,a$ (and the pair sum is zero for $m=1$).
Its total integral has a deterministic bound. Namely,
\[
 \int_0^{U_i}\delta_i(u)\,\dd u
 =\frac2\beta(1-e^{-\beta U_i/2})\le\frac2\beta,
 \qquad
 (\delta_i\delta_j)^{-M}\le e^{M\beta(u_i+u_j)}.
\]
Equation~\eqref{mp:interaction-integral}, with $c=M\beta$, yields
\begin{equation}
 \int k_t\,\dd t\le
 K_{m,a}:=\frac{2mC_1}{\beta}
             +C_3\binom m2 I(a,M\beta)<\infty.
 \label{mp:discount}
\end{equation}
It follows that
$Q_t\exp(-\int_0^t k_v\,\dd v)$ is a nonnegative local
supermartingale on every nonsingular interval, including its downward
activation and completion jumps.

\begin{proof}[Completion of the proof of Proposition~\ref{mp:estimate}]
To justify stopping the product, monitor \emph{all unfinished points},
including those not yet active. Stop if any of their angles leaves
$(1/N,\pi-1/N)$, if any of their heights $Y_i$ falls to $1/N$, or at
physical time $N$. Denote this cutoff by $\zeta_N$, and also stop when
all targets have been completed. No cutoff is imposed on a completed
point. At these stops every stochastic integral used above is
legitimate; a further standard bounded localizer, followed by Fatou,
gives expectation at most the initial value, which is at most $1$.
There is no need to continue a factor through the singular lifetime of
an unsuccessful unfinished point.

Put $T_{\rm all}=\max_iT_i$. On $\{T_{\rm all}<\zeta_N\}$, the final
discounted product is at least
\[
 e^{-K_{m,a}}3^{-m}\exp\left(\frac p2\sum_iU_i\right).
\]
Consequently
\[
 \PP(T_{\rm all}<\zeta_N)
 \le 3^m e^{K_{m,a}}
                   \exp\left(-\frac p2\sum_iU_i\right).
\]
Every successful path belongs to these events for all sufficiently
large $N$. Indeed, each point remains alive through its own finite
$T_i$, so its continuous height and its distance in angle from the
endpoints have strictly positive minima on $[0,T_i]$. There are only
finitely many points, and later swallowing of a finished point is
irrelevant. Thus cutoff exhaustion covers precisely all the success
paths needed for the upper bound. Since $p/(4a)=\alpha$, passage to
the limit gives \eqref{mp:hit-bound} with
$C_{m,a}=3^m e^{K_{m,a}}$ when all $e_i<b_i$.

For equality targets, or targets larger than their caps, retain only
the strict constraints $e_i<b_i$. The original event is contained in
the retained event, and all separation assumptions still hold.
Taking the maximum of the constants for at most $m$ retained points,
including the empty set, proves \eqref{mp:capped-hit-bound} and the
full stated proposition.
\end{proof}

\paragraph{Comparison with the Euclidean estimate.}
To deduce the proposition from Rezaei--Zhan, put $f=J^{-1}$,
$z_i=f(x_i)$, $\lambda_i=|f'(x_i)|$, and $z_0=0$.
Koebe's Quarter Theorem applied to $f$ on $B(x_i,b_i/4)\subset D$
gives
\begin{align*}
 |z_i-z_j|&\ge\lambda_i b_i/16
                   &&(1\le j\le m,\ j\ne i),\\
 \operatorname{Im}z_i
 &=\lambda_iR_D(x_i)/2\ge\lambda_i b_i/2.
\end{align*}
Indeed, $f(B(x_i,b_i/4))$ contains the disk with center $z_i$ and
radius $\lambda_i b_i/16$, whereas every other $x_j$ lies outside
$B(x_i,b_i/4)$. Retain only targets with $e_i<b_i/32$ and relabel
the retained points; if none remain, the trivial probability bound
suffices. Their preceding-point distances satisfy
$l_i=\min_{0\le j<i}|z_i-z_j|\ge\lambda_i b_i/16$.
At a successful hit, radius covariance gives
$R_{H_t}(z_i)=\lambda_i e_i$. Koebe's distance
bound then puts the past half-plane trace within
$r_i=\lambda_i e_i$ of $z_i$: the real line is farther away.
Theorem~1.1 and equation~(1.5) of \cite{RezaeiZhan2017} give
\[
 \PP\{\dist(z_i,\gamma)\le r_i\text{ for all retained }i\}
 \le C_{m,a}\prod_{i\text{ retained}}(r_i/l_i)^\alpha
 \le C_{m,a}16^{m\alpha}
             \prod_{i\text{ retained}}(e_i/b_i)^\alpha.
\]
Each omitted target contributes at least $32^{-\alpha}$ to the capped
product, so it can
be restored by enlarging the constant. This proves the stated
domain-uniform consequence without a boundary extension of $J$.

\subsection{Integrating the collision singularities}

\begin{lemma}[Capped distance kernels]\label{mp:kernel-integral}
Let $0<\alpha<1$, $b>0$, and let $K\subset\mathbb C$ be a bounded Borel set.
For distinct $x_1,\ldots,x_m$ set
\[
 b_i=\min\bigl(b,\min_{j\ne i}|x_i-x_j|\bigr),\qquad
 F_b(x_1,\ldots,x_m)=\prod_{i=1}^m b_i^{-\alpha},
\]
where the inner minimum is omitted when $m=1$. Then
\begin{equation}
 \int_{K^m}F_b(x_1,\ldots,x_m)\prod_i\dd A(x_i)<\infty.
 \label{mp:all-variable-integral}
\end{equation}
The same assertion holds with one point fixed: for bounded $K_0$,
\begin{equation}
 \sup_{x_0\in K_0}\int_{K^m}
 F_b(x_0,x_1,\ldots,x_m)\prod_{i=1}^m\dd A(x_i)<\infty.
 \label{mp:fixed-center-integral}
\end{equation}
Values on diagonals may be set to infinity; they do not affect these
integrals.
\end{lemma}

\begin{proof}
For each factor,
$b_i^{-\alpha}\le b^{-\alpha}+\sum_{j\ne i}|x_i-x_j|^{-\alpha}$.
Expand their product. Each resulting term corresponds to a directed
graph with at most one outgoing edge per vertex, no self-edge, and a
factor $|x_i-x_j|^{-\alpha}$ for each edge $i\to j$.
Integrate a vertex with no incoming edge: if it has one outgoing
edge its integral is bounded uniformly in the other endpoint because
$\alpha<2$; if isolated it contributes only the area of a containing
box. Repeating removes all trees attached to directed cycles and all
components without cycles.

For a remaining cycle of length at least three, one vertex has two
incident factors. Their integral is uniformly bounded by
\[
 \int_K|x-y|^{-\alpha}|x-z|^{-\alpha}\,\dd A(x)
 \le\left(\int_K|x-y|^{-2\alpha}\,\dd A(x)\right)^{1/2}
      \left(\int_K|x-z|^{-2\alpha}\,\dd A(x)\right)^{1/2}
 \le C_K,
\]
since $2\alpha<2$. Removing this vertex leaves a chain, whose
endpoints can be integrated successively. A two-cycle contributes
$|x-y|^{-2\alpha}$, so the same integrability applies directly.
There are finitely many graph terms, proving
\eqref{mp:all-variable-integral}.

For the fixed point, choose $R<\infty$ such that
$K-x_0\subset B(0,R)$ for every $x_0\in K_0$. The kernel $F_b$ is
invariant under simultaneous translation, with the cap $b$ unchanged.
Thus the integral in \eqref{mp:fixed-center-integral} is bounded by
\[
 \int_{B(0,R)^m}F_b(0,y_1,\ldots,y_m)\prod_i\dd A(y_i).
\]
Average a simultaneous translation by $v\in B(0,1)$. A change of
variables bounds the last display by
\[
 \frac1{|B(0,1)|}
 \int_{B(0,R+1)^{m+1}}F_b(v,x_1,\ldots,x_m)
                       \,\dd A(v)\prod_i\dd A(x_i),
\]
which is finite by the already proved assertion. This is translation
invariance of a deterministic distance kernel; no translation
invariance of a chordal SLE law is required.
\end{proof}

We record precisely the moment consequence used in the next sections.
For $l>0$ let
\[
 A_l(\dd x)=G_{\tau_l(x)}(x)
                       \ind_{\{\tau_l(x)<\infty\}}\,\dd A(x).
\]

\begin{corollary}[Integrated level weights]\label{mp:level-moments}
Let $f\ge0$ be bounded and supported in a compact $K\Subset D$.
Then, for every integer $m\ge1$,
\begin{equation}
 \sup_{l>0}\EE[A_l(f)^m]<\infty.
 \label{mp:ordinary-moments}
\end{equation}
Fix also $x_0\in D$ and $0<e_0\le R_D(x_0)$, and let $E_0$ be the
event that $x_0$ reaches radius $e_0$ while alive. Then
\begin{equation}
 \sup_{l>0}\EE[\ind_{E_0} A_l(f)^m]\le C e_0^\alpha.
 \label{mp:center-moments}
\end{equation}
Here $C$ can be chosen using only $m,a$, the size and support of $f$,
the bounded location of $x_0$, and a positive lower bound for
$R_D$ on $K\cup\{x_0\}$. In particular, in a fixed disk these bounds
are uniform over all initial boundary angles.
\end{corollary}

\begin{proof}
Choose a fixed cap $b>0$ below the indicated initial radii, and cap
it further by nearest-neighbor distances for each integration tuple.
At point $x_i$, the target is $e_i=R_D(x_i)e^{-2al}$, and
$G_{\tau_l(x_i)}(x_i)\le e_i^{-\alpha}$ on success. Tonelli and
\eqref{mp:capped-hit-bound} therefore bound the $m$th moment by
\[
 C_{m,a}\int_{K^m}\prod_{i=1}^m f(x_i)b_i^{-\alpha}
                                                 \prod_i\dd A(x_i),
\]
because $e_i^{-\alpha}\min\{1,(e_i/b_i)^\alpha\}\le b_i^{-\alpha}$.
This proves \eqref{mp:ordinary-moments} by
Lemma~\ref{mp:kernel-integral}. For the additional fixed center, apply
the same argument to $m+1$ points and use
\[
 \min\{1,(e_0/b_0)^\alpha\}\le e_0^\alpha b_0^{-\alpha}.
\]
The fixed-center integral proves \eqref{mp:center-moments}, even
when the center's target exceeds its distance cap. No relative order
of the center hit and the other hits is imposed. Finally, in $\DD$
one has $R_D(x)=1-|x|^2$, independently of its boundary marks. The
proof was carried out directly in the physical domain, so a
degenerating half-plane identification introduces no additional
constant.
\end{proof}
